\documentclass[11pt]{report}
\usepackage[utf8]{inputenc}
\usepackage{amsmath}
\usepackage{amsfonts}
\usepackage{amssymb}
\usepackage{fancyhdr}
\usepackage{graphicx}
\usepackage{lipsum}
\usepackage{geometry}
\usepackage{lastpage}
\usepackage{indentfirst}
\usepackage{titlesec}
\usepackage{blindtext, color}
\usepackage{hyperref}
\usepackage{nameref}
\usepackage{tabularx}
\usepackage{glossaries}
\usepackage{xcolor}
\usepackage[toc,page,title,header]{appendix}
\usepackage{enumerate}
\usepackage{caption}
\usepackage[normalem]{ulem}

%-----Couleur des liens cliquables-----%
\definecolor{bleu marine}{HTML}{1E3F5A}
\hypersetup{
    colorlinks=true,
    linkcolor=gray,
    filecolor=magenta,
    urlcolor=bleu marine,
    citecolor = gray,
}
%---------------------------------------%

%-----Mise en forme des légendes-----%
\DeclareCaptionFont{gras}{\bfseries}
\DeclareCaptionLabelFormat{test}{\small\uline{#1~#2}}
\DeclareCaptionTextFormat{test}{\small\uline{#1}}
\DeclareCaptionLabelSeparator{test}{\small\uline{:~}}

\captionsetup{
format=hang,labelformat=test,textformat=test,labelseparator=test, labelfont = gras}
%------------------------------------%

%-----Mise en forme des titres-----%
\definecolor{gray75}{gray}{0.75}
\newcommand{\hsp}{\hspace{20pt}}
\titleformat{\chapter}[hang]{\huge\bfseries}{\thechapter\hsp\textcolor{gray75}{|}\hsp}{0pt}{\bfseries}
\titlespacing*{\chapter}{0pt}{-50pt}{20pt}
\titleformat{\subsection}[hang]{}{\hspace{1em}\bfseries\thesubsection}{1em}{\bfseries}
\titleformat{\subsubsection}[hang]{}{\hspace{2em}\bfseries\thesubsubsection}{1em}{\bfseries}
\renewcommand{\thesubsubsection}{\alph{subsubsection})}
%----------------------------------%

%-----Redéfinition des noms des parties-----%
\renewcommand{\contentsname}{Sommaire}
\renewcommand{\tablename}{Tableau}
\renewcommand{\chaptername}{Chapitre}
\renewcommand\listfigurename{Liste des Figures}
\renewcommand\listtablename{Liste des Tableaux}
\renewcommand{\appendixtocname}{Annexes}
\renewcommand{\appendixpagename}{Annexes}
\renewcommand{\bibname}{Bibliographie}
%-------------------------------------------%

%-----Création nouveau type de colonne-----%
\newcolumntype{M}[1]{>{\centering\arraybackslash}m{#1}}
%------------------------------------------%

%-----Mise en page globale-----%
\graphicspath{{photo/}}
\pagestyle{fancy}

\fancypagestyle{plain}{
    %Chapter appelle le style plain donc obligé de faire ça%
    \fancyhead[L]{\leftmark}
    \fancyhead[R]{\includegraphics[scale = 0.3]{logo phelma.png}}
    \renewcommand\footrulewidth{0.5pt}
    \fancyfoot[C]{\textbf{Page \thepage/\pageref{LastPage}}}
    \setlength{\headsep}{45pt}
}

\fancyhead[L]{\leftmark}
\fancyhead[R]{\includegraphics[scale = 0.05]{logo phelma.png}}
\renewcommand\footrulewidth{0.5pt}
\fancyfoot[C]{\textbf{Page \thepage/\pageref{LastPage}}}
\setlength{\headheight}{46pt}
\setlength{\headsep}{15pt}
\setlength{\topmargin}{-26pt}
\setlength{\textwidth}{15cm}
\setcounter{secnumdepth}{4}
\setcounter{tocdepth}{4}
%-------------------------------%

%-----GLOSSAIRE-----%

%--------------------%


\begin{document}

%-----Première page-----%
\begin{titlepage}
    \hfill
    \begin{center}
        \begin{minipage}{1\linewidth}
        \begin{minipage}{.5\linewidth}
              \begin{flushleft} %aligner à gauche
             \includegraphics[scale=0.5]{logo phelma.png} %% Logo haut gauche
            \end{flushleft}
        \end{minipage}
        \begin{minipage}{.5\linewidth}
              \begin{flushright}    
         %	\includegraphics[scale=.5]{logo2} %% Logo haut droite
            \end{flushright}
        \end{minipage}
        \end{minipage}
        
        \vspace*{\stretch{1}}
        
        \textsc{\LARGE Dialogue entre dynamique moléculaire et milieux continu: définition d'un modèle cohésif fondé sur des informations à l'échelle atomique}

        \bigskip
        \bigskip
        \bigskip
        
        \hrule height 2px \bigskip
        { \huge \bfseries - RAPPORT DE STAGE 2A - \\[0.4cm] }
        
        \hrule height 2px 

        \bigskip
        {\large Année universitaire 2021-2022} \\[0.7cm] 
        
        \begin{figure}[!h]
            \centering
            \includegraphics[scale=1.2]{logo_SIMaP_801.png}
        \end{figure}
              \vfill 
        
        \begin{minipage}[d]{1\linewidth}
            \begin{flushright} \large
            \textbf{Maître de Stage}\\
                \textbf{Noel JAKSE}\\
                0102030405\\
                noel.jakse@grenoble-inp.org
            \end{flushright}
        \end{minipage}
        
        \bigskip
             
        \begin{minipage}[d]{1\linewidth}
              \begin{flushleft} \large
              \textbf{Thibault MROZ}\\
                \textbf{2A - SIM}\\
              \end{flushleft}
         \end{minipage}
          \end{center}        
\end{titlepage}
%-----------------------%

\tableofcontents

\setcounter{chapter}{0}
    \chapter*{Remerciements}\chaptermark{Remerciements}
    \addcontentsline{toc}{chapter}{Remerciements} %car non numéroté%

    % \chapter*{Abstract et Résumé}\chaptermark{Abstract et Résumé}
    % \addcontentsline{toc}{chapter}{Abstract et Résumé}


%\printglossary[title={Glossaire}, toctitle={Glossaire}]%

    %\addcontentsline{toc}{chapter}{Glossaire}%

	\listoffigures
	\addcontentsline{toc}{chapter}{Liste des Figures}

    \listoftables
    \addcontentsline{toc}{chapter}{Liste des Tableaux}

\chapter{Introduction}
    % Blabla Introduction

\section{Contexte} % Expliquer dans quel contexte s'inscrit le stage

    Les avancées technologiques de l'Intelligence Artificielle et du Machine Learning ont permis l'aboutissement de beaucoup d'outils facilitant la vie des utilisateurs et de méthodes pour faire avancer la recherche. Phénomène très récent dans le monde de la recherche, le Machine Learning permet de gagner beaucoup de temps de calculs pour réaliser des simulations à grande échelle mais également réaliser quelques prédictions. 
    
    Il y a plusieurs branches de Machine Learning. La plus classique est celle où on fourni à un programme beaucoup de données et il apprend à partir de ces données. Le programme pourra alors fournir une prédiction en fonction des paramètres d'entrée. Cependant, en faisant ça, on perd la physique derrière (si les données sont des simulations ou expériences physiques). Une autre branche est de faire du tri sur les données en entrée en fonction de ce qui est plus probable d'arriver. Le sens physique est alors conservé mais le programme ne pourra fournir qu'une estmiation de ce qui pourrait arriver. 

\section{Problématique du Stage}

    Ce stage s'inscrit dans le cadre d'une analyse multi-échelles de la rupture et plus précisément en menant un diaglogue en dynamique moléculaire et description en milieu continu. Plus précisément, il s'agit d'identifier un modèle de zone cohésive, représentant à l'échelle continu le mécanisme de rupture au travers d'une relation "vecteur-contrainte" - "ouverture". Ce modèle sera identifiable suite à des calculs en dynamique moléculaire qui produisent les expérimentations numériques. \medskip

    Une partie importante du travail est de mener des simulations de Dynamique Moléculaire sur un cristal de Silicium (Si) pour lequel la rupture a lieu par clivage. Cependant, le Si possédant des propriétés élastiques anisotropes, il est attendu que les propriétés de rupture le soient également. Dès lors, des simulations pour différentes orientations entre plan de la fissure / plans de symétrie critallins seront également à mener. Une approche systématique peut être menée. Néanmoins, la méthodologie associant Machine Learning et Dynamique Moléculaire est à exploiter afin de gagner en temps de calculs. \medskip
    
    Une fois que le modèle cohésif est identifié, il est ensuite possible d'étudier et de prédire les intercations entre fissuration et microstructure (dans un polycristal par exemple), ainsi qu'entre fissure et cavité. \medskip

    Il s'agit d'un projet "100\% numérique" et porte un fort intérêt pour les méthodes de simulations ainsi que le Machine Learning.

\section{Présentation du Laboratoire}

    \subsection{Historique}

        Le Laboratoire de Science et Ingénierie des Matériaux et Procédés résulte de la fusion de trois unités au 1er Janvier 2007. C'est une UMR, Unité Mixte de Recherche : CNRS, Grenoble-INP, et IESA. Il rassemble en moyenne 220 personnes dont 56 chercheurs et enseignants-chercheurs, 37 ingénieurs, techniciens et administratifs, 60 doctorants, les post-doctorants, invités et stagiaires. 

    \subsection{Groupes de Recherche}

        Le Laboratoire s'appuie sur quatre groupe de recherche qui pérennisent les sciences de base en physique et physico-chimie, thermodynamique et cinétique, mécanique des solides et des fluides : 

        \begin{enumerate}[\hspace{3em}$\bullet$]
            \item \textbf{EPM : }Élaboration par Procédés Magnétiques
            \item \textbf{GPM2 : }Génie Physique et Mécanique des Matériaux
            \item \textbf{PM : }Physique du Métal
            \item \textbf{TOP : }Thermodynamique, modélisation, Optimisation des Procédés
        \end{enumerate}
    
        Ce stage se place entre deux divisions (PM et TOP) dans une petite équipe constituée de : 

        \begin{enumerate}[\hspace{3em}$\bullet$]
            \item \textbf{Noel JAKSE : }Enseignant-Chercheur au groupe de recherche TOP, Maître et Tuteur de Stage
            \item \textbf{Rafael ESTEVEZ : }Chercheur, Co-tuteur de Stage
            \item \textbf{Thibault MROZ : }Stagiaire Assistant Ingénieur
        \end{enumerate}

        Le Groupe de Recherche TOP se concentre sur l'élaboration des matériaux, les phéno-
        mènes thermodynamiques (stabilité et charactérisation) et la modéliation atomistique, thermo-
        dynamique, cinétique et des réacteurs. Cela a des applications dans les domaines des films minces, des alliages métalliques complexes et des matériaux fonctionnels.

        Le Groupe de Recherche PM se concentre sur la métallurgie des métaux : structure atomique, propriétés mécaniques et physiques ainsi que l'oxydation. Cela a des applications dans les domaines des matériaux pour l'énergie et pour la micro-électronique mais aussi pour les matériaux structuraux. 

        Le stage s'inscrit dans le domaine de la modélisation atomistique et de la structure atomique. 

    \subsection{Présentation et Déroulement du Stage}

\chapter{Gestion de Projet}


    \section{Répartition des rôles}


\chapter{Mise en contexte}

    \section{Contexte d'utilisation}


\chapter[Étude d'hydroliennes en milieu artificiel]{Étude des caractéristiques des prototypes d'hydroliennes en milieu artificiel}

    \section{Conception et réalisation du système expérimental}

\chapter[Implantation des hydroliennes sur Grenoble]{Étude du potentiel d'implantation des hydroliennes sur Grenoble}



\chapter*{Conclusion et Perspectives}\chaptermark{Conclusion et Perspectives}
\addcontentsline{toc}{chapter}{Conclusion et Perspectives}


\begin{thebibliography}{9}%Voir sur internet pour le 0 
%Revoir biblio%
    \addcontentsline{toc}{chapter}{Bibliographie}
\end{thebibliography} 

%-----ANNEXES-----%
\begin{appendices}


\fancyhead[L]{ANNEXES}
\fancyhead[R]{\includegraphics[scale = 0.05]{logo phelma.png}}
\renewcommand\footrulewidth{0.5pt}
\fancyfoot[C]{\textbf{Page \thepage/\pageref{LastPage}}}
\setlength{\headheight}{46pt}
\setlength{\headsep}{15pt}
\setlength{\topmargin}{-26pt}
\setlength{\textwidth}{15cm}
\setcounter{secnumdepth}{4}
\setcounter{tocdepth}{4}


\renewcommand{\thesection}{\Alph{section})}

\section{Annexe 1}
\label{section:annexe_1}

\end{appendices}
%--------------------------------------------------%


\end{document}
